\subsection{A mixed chart with a complete formal proof}\label{sec:mixedformal}

One of the previously analytic exclusions admits a direct formal proof.
It is the whole family in the next proposition, not a generic choice
of its coefficients.

\begin{proposition}\label{prop:mixedformal}
For every $a\in\C$ and every $G\in\C[Z,T]$ of total degree at most
three, the polynomial
\[
 P=W(a+T+Z^2)+G(Z,T)
\]
has a nonzero positive pure Gaussian moment.
\end{proposition}
\begin{proof}
We first construct the exact formal seed functional, then a full frame
at infinity. This proof does not identify a divergent Gaussian series
with an analytic integral.

\emph{The formal source.}
The circular fixed point is $g=t(a+T+g^2)$, with determinant
$D=1-2tg$. Make the formal Gaussian change
\[
 T=Y-t^2(a+Y)^2,\qquad z=t(a+Y).
\]
It has $g=z$ and $dT/dY=D$. For any change $T(Y)=Y+O(t)$ with
polynomial formal coefficients, the identity
\[
 \E_T[F(T)]=\E_Y\!\left[F(T(Y))T'(Y)
             e^{(Y^2-T(Y)^2)/2}\right]
\]
holds coefficientwise. To prove it, use the homotopy
$T_u=Y+u(T-Y)$: its derivative is the Gaussian expectation of
$(\partial_Y-Y)\{F(T_u)\dot T_u e^{(Y^2-T_u^2)/2}\}$, which is zero.
Thus the determinant cancels and gives
\[
 \Phi_1=\E_Y[e^\Psi],\qquad
 \Psi=tG(z,Y-z^2)+Yz^2-z^4/2.
\]

Define
\[
 F=z/t-a-z^2,\qquad \varphi(t,z)=tG(z,F)-F^2/2,
\]
and let
\[
 M=K[z]/(\partial_z+\varphi_z)K[z],\qquad
 \nabla[q]=[q_t+(\varphi_t-1/t)q],\qquad K=\C(t).
\]
These are vector-space quotients. The formal map
$J(q)=\E_Y[q(t(a+Y))e^\Psi]$ annihilates
$(\partial_z+\varphi_z)K[z]$ by Gaussian integration by parts.
Differentiating at fixed $Y$ differs from applying $\nabla$ by
$J((\partial_z+\varphi_z)(zq/t))=0$.
So $J$ is horizontal and $J(1)=\Phi_1$. The term $-1/t$ is the
normalization $dz/t=dY$; removing it changes the problem.
Leading-degree reduction gives $\dim_K M=\deg_z\varphi-1$.

\emph{The exhaustive degree split.}
Put
\[
 V(z)=G(z,-a-z^2),\quad L(z)=G_T(z,-a-z^2),\quad
 B_1=z^2G_{TT}(z,-a-z^2)/2+z(a+z^2).
\]
If $s=[T^3]G$, direct Taylor expansion gives
\begin{equation}\label{eq:mixedphi}
 \varphi=tV+B_0+B_1/t+B_2/t^2,
 \quad B_0=zL-(a+z^2)^2/2,\quad B_2=sz^3-z^2/2.
\end{equation}
To check every degree drop, write $v_1=[ZT^2]G$,
$v_0=[T^2]G$ and $u_2=[Z^2T]G$. Then
\[
 [z^6]V=-s,\qquad [z^5]V=v_1,\qquad
 [z^4]V=v_0-u_2-3as.
\]
If $n=\deg V\geq4$, the degree of $\varphi$ is $n$ and each $B_j$
has degree at most $n$. If $n\leq3$, or $V$ is constant, these three
identities force $s=v_1=0$ and $u_2=v_0$. Hence
$B_0=-z^4/2+O(z^3)$, $\deg B_1\leq3$ and $B_2=-z^2/2$.
In that case $\deg_z\varphi=4$. There is no omitted higher-degree
term tending to zero at infinity.

\emph{Finite characters.}
At a critical point $p$ of the fixed polynomial $V$, write locally
$V=\alpha+y^k$, $k\geq2$, in a formal invertible coordinate.
The $k-1$ normalized power functionals are indexed by $1\leq j<k$:
their value on $y^b$ is zero unless $b=j-1+qk$, $q\geq0$;
in the latter case its value, up to a fixed nonzero normalization, is
\[
 e^{\alpha t}t^{-j/k}(-1)^q(j/k)_q t^{-q}.
\]
The recursion $I(y^{b+k})=-(b+1)I(y^b)/(kt)$ proves spatial
integration by parts. Apply these functionals to the Taylor expansion
of the polynomial amplitude times $e^{B_0+B_1/t+B_2/t^2}dz/dy$.
Each formal coefficient uses finitely many Taylor terms. The recursion
also proves parameter horizontality. Including the factor $1/t$, the
resulting functionals have factors
\[
 e^{\alpha t}t^{-1-j/k}\C[[t^{-1/k}]].
\]
They retain multiple critical points; splitting them into Morse points
is unnecessary. Their total number is $n-1$ when $V$ is nonconstant.

\emph{The case $n\geq4$.}
The finite characters already number $n-1=\dim M$.
Confluent polynomial interpolation at the critical points makes their
leading evaluation matrix invertible. Indeed their leading evaluations
are the jets of orders $0,\ldots,k-2$, multiplied by an analytic unit
and an invertible local coordinate change. The sum of the jet lengths
is $n-1$. This proves completeness, rather than assuming that the
constant amplitude is cyclic. Each zero-exponential seed value has
strictly positive order at infinity.

\emph{The cases $1\leq n\leq3$.}
There are $n-1$ finite characters. Put $r=t^{-1/(4-n)}$ and $z=\zeta/r$.
The leading phase in \eqref{eq:mixedphi} is
\[
 r^{-4}H(\zeta),\qquad H(\zeta)=V_n\zeta^n-\zeta^4/2.
\]
All other terms have strictly higher $r$-order. Its $4-n$ nonzero
critical points satisfy
\[
 \zeta^{4-n}=nV_n/2,\quad
 H(\zeta)=(1-n/4)V_n\zeta^n\ne0,\quad
 H''(\zeta)=2(n-4)\zeta^2\ne0.
\]
The formal implicit function theorem gives the perturbed centers
$\zeta_i(r)$. Substitution $\zeta=\zeta_i(r)+r^2U$ gives a
nondegenerate fixed quadratic form in $U$ and a perturbation divisible
by $r$. Formal Gaussian contraction gives one horizontal functional
per center. Its exponential part $q_i$ has the nonzero leading term
$r^{-4}H(\zeta_i(0))$, and its volume factor, including $1/t$, is
$r^{5-n}$. Thus all escaping factors have nonzero exponential part.
Their physical $z$ fluctuations have positive order $O(r)$.

Here is an explicit completeness check. Let $\mu_p=k_p-1$,
$F=\sum_p\mu_p=n-1$, $D(X)=\prod_p(X-p)^{\mu_p}$ and
$E(X,r)=\prod_i(X-z_i(r))$, where $z_i=\zeta_i(r)/r$.
For each finite point choose its confluent Hermite polynomials
$H_{p,j}$ modulo $D$, with $0\leq j<\mu_p$. Use the following
three polynomial amplitudes:
\[
 H_{p,j}(z)\left(\frac{E(z,r)}{E(p,r)}\right)^4,
 \qquad
 \frac{D(z)}{D(z_i)}\prod_{\ell\ne i}
              \frac{z-z_\ell}{z_i-z_\ell}.
\]
Terms involving finite points are absent when $F=0$.
At a bounded point $E(z,r)/E(p,r)=1+O(r)$, so the finite block is
the invertible jet matrix just described. Near an escape, this ratio
is $O(r^2)$, while $H_{p,j}$ grows at most as $r^{-(F-1)}$.
Thus the finite rows vanish in the normalized escape columns.
The escape rows have nonzero diagonal Gaussian evaluations and
vanishing cross terms, since all escape separations are of order
$r^{-1}$ and all fluctuations are $O(r)$. Their finite evaluations
vanish to positive order as well, since $D(z_i)^{-1}=O(r^F)$
when finite points exist. The full matrix is invertible. The count
$(n-1)+(4-n)=3$ equals the complete de Rham rank.

\emph{The constant case.}
If $V=k$, write $\varphi=tk+\varphi_0(z)+O(t^{-1})$, where
$\varphi_0$ has degree four and leading coefficient $-1/2$.
Its twisted quotient has dimension three. For any basis $\ell_i$
of the dual of this fixed quotient, define formally
\[
 q\longmapsto e^{tk}t^{-1}
     \ell_i\!\left(e^{\varphi-tk-\varphi_0}q\right).
\]
At each order in $t^{-1}$ the expression in parentheses is a
polynomial. Product differentiation proves both spatial integration
by parts and horizontality. The leading evaluation matrix on a fixed
quotient basis is invertible, even if $\varphi_0$ has degenerate critical
points. If $k=0$, all seed values have a factor $t^{-1}$.

In each regime we have constructed a full formal frame. The only
zero-exponential seed values have positive order after a finite
ramification. Lemma~\ref{lem:positiveformal} excludes the quotient
with seed value one that purity would supply. This proves the proposition.
\end{proof}
